\section{Introduction and Physical Model}
\subsection{Robertson reaction system}
The Robertson problem is a standard stiff reaction-rate benchmark~\cite{robertson}. With $y=(y_1,y_2,y_3)^T=([A],[B],[C])^T$,
\begin{align}
y_1'&=-0.04y_1+10^4y_2y_3,\qquad
y_2'=0.04y_1-10^4y_2y_3-3\times10^7y_2^2,\label{eq:model}\\
y_3'&=3\times10^7y_2^2,
\end{align}
with $y(0)=(1,0,0)^T$ on $0\le t\le40$. The variables are normalised concentrations of $A,B,C$.

Adding the equations gives $(1,1,1)^Tf(y)=0$, hence
\[
y_1(t)+y_2(t)+y_3(t)=1.
\]
Species $B$ is a short-lived intermediate. On the prescribed logarithmic output grid, the sampled maximum is $y_2\approx3.6487\times10^{-5}$ at $t\approx4.42\times10^{-3}\,\mathrm{s}$, while $y_2(40)=9.185535\times10^{-6}$ and $A$ and $C$ evolve over tens of seconds.

\begin{figure}[H]
\centering\includegraphics[width=\textwidth]{F2_reference_solution.pdf}
\caption{Tight-tolerance reference trajectory. Panels separate the slow species, the small intermediate $B$, and the invariant defect. Logarithmic time resolves the early transient; the invariant remains near floating-point round-off.}
\label{fig:reference}
\end{figure}

\subsection{Why the problem remains stiff after the visible transient}
The Robertson system remains stiff after the visible spike in $B$ has passed. A large negative Jacobian eigenvalue can persist while the concentration curve varies slowly, so the explicit step restriction follows from the decay modes rather than the apparent shape of the trajectory. We quantify this separation through the eigenvalues and characteristic times $1/|\operatorname{Re}\lambda|$.

\begin{figure}[H]
\centering\includegraphics[width=\textwidth]{F13_two_clocks.pdf}
\caption{Concentrations and characteristic times. The slow chemistry determines the 40-second horizon, whereas the persistent fast clock restricts explicit steps. The clock ratio grows from about $1.7$ to $1.58\times10^5$.}
\label{fig:clocks}
\end{figure}

\subsection{Focus and report structure}
This report studies the Robertson problem through reference-solution validation, stiffness analysis, cost--accuracy comparison, and experiments on Newton and adaptive tolerances. The numerical methods are introduced first, followed by stability analysis and numerical experiments.

\subsection{Scope of analysis}
The report applies standard numerical methods to the stiff Robertson system. For implicit Euler, the Newton-tolerance experiment examines how predictor acceptance affects the computed update and trajectory error. A loose nonlinear tolerance can accept the explicit-Euler predictor without correction and change the update realised by the nominally implicit method.

Numerical reliability is checked through conservation, non-negativity, full-state error, and the endpoint error in $y_2(40)$. Each diagnostic captures a different type of numerical failure. The explicit-Euler counterexample preserves the invariant near machine precision while producing negative concentrations and a large trajectory error.

The adaptive study shows how the error norm changes acceptance behaviour, global error, and physical admissibility. Robertson concentrations occupy very different scales, so a small absolute error can still be large relative to the intermediate species. The comparison uses an unscaled absolute infinity norm and component-scaled mixed control over the tested parameter range.

The adaptive study is extended by converting the frozen-Jacobian stability estimate into a local cap on the RK4 step proposal and comparing it with the standard mixed-norm controller.

\begin{table}[H]
\centering
\footnotesize
\caption{Summary of analyses performed. The additional analyses use standard numerical methods rather than new algorithms.}
\label{tab:extensions}
\begin{tabular}{@{}>{\raggedright\arraybackslash}p{0.18\textwidth}>{\raggedright\arraybackslash}p{0.34\textwidth}>{\raggedright\arraybackslash}p{0.38\textwidth}@{}}
\toprule
Required element&Additional analysis&Numerical purpose\\
\midrule
Stability analysis&Frozen-Jacobian prediction combined with empirical step-size bracketing and separate tracking of negativity and divergence&Tests whether local linear stability predicts the relevant scale rather than an exact nonlinear boundary\\
Implicit Euler&Predictor-acceptance and Newton-residual analysis&Shows how nonlinear stopping tolerance can change the realised update and expose explicit-like behaviour\\
Convergence study&Pre-asymptotic interpretation using $E(h)=C_ph^p+C_{p+1}h^{p+1}+\cdots$&Explains why finite-range fitted slopes need not equal the formal order exactly\\
Adaptive integration&Absolute-norm versus component-scaled error-control comparison&Tests whether an apparently acceptable local error can still fail to protect the small intermediate species\\
\bottomrule
\end{tabular}
\end{table}

The experiments compare stability, accuracy, and non-negativity because these criteria can give different conclusions for the same numerical run.

The large separation of chemical time scales motivates the comparison of explicit and implicit time integrators. We therefore introduce their accuracy and linear-stability properties before applying trajectory-dependent Jacobian analysis to the Robertson system.
